\chapter{اللقاء الثالث: تابع مقدمة الطبري وقضية الإعجاز}

\section{مقدمة عن نعمة البيان وغاية الآيات}
السلام عليكم ورحمة الله وبركاته. الحمد لله رب العالمين، والصلاة والسلام على نبينا محمد وعلى آله وصحبه أجمعين. هذا هو اللقاء الثالث من مجالسنا في دراسة «تفسير الطبري». 

توقفنا عند مقطع مهم من مقدمة الإمام \rawi{الطبري} يتكلم فيه عن اتفاق معاني آي القرآن ومعاني منطق من نزل بلسانه، ومن وجوه البيان.

\isnad{قال الإمام} \rawi{الطبري} \isnad{رحمه الله:}
«القول في البيان عن اتفاق معاني منطق مَن نزل بلسانه، ومن وجه البيان، والدلالة على أن ذلك من الله هو الحكمة البالغة، مع الإبانة عن فضل المعنى الذي باين به القرآن سائر الكلام. اعلم أن من عظيم نعم الله على عباده وجسيم منّته على خلقه، ما منحه من فضل البيان الذي به عن ضمائر صدورهم يُبينون، وبه على عزائم نفوسهم يدلون...».

\begin{fawaid}[كيف تقرأ مقدمات الكتب التراثية؟]
حينما تقرأ مقدمة لكتاب ما، هناك ثلاثة أمور رئيسية يجب أن تعتني بها وتستخرجها:
\begin{enumerate}
    \item \textbf{حصر موضوعات المقدمة:} ما هي القضايا الكلية التي تطرق إليها المؤلف؟
    \item \textbf{حصر المسائل المندرجة:} ما هي المسائل الدقيقة تحت كل موضوع؟
    \item \textbf{خلاصات المقدمة:} ما هي النتائج النهائية التي أراد المصنف أن يصل إليها؟
\end{enumerate}
\end{fawaid}

\section{آيات الأنبياء ومدخل الطبري إلى الكلام على القرآن}
يُريد \rawi{الطبري} أن يمهد للحديث عن القرآن بوصفه آية النبي صلى الله عليه وسلم العظمى، فقدم لذلك بالكلام عن نعمة البيان؛ لأن البيان هو أشرف ما تفاضل به العرب في باب اللسان، ومن هنا كان ظهور عجزهم أمام القرآن أدخل في الحجة وأبين في البرهان.

ثم انتقل الشيخ من هذا الموضع إلى مسألة دقيقة تتعلق بما يسمى اصطلاحاً بـ "المعجزة". فقد جعل عدد من المتكلمين آيات الأنبياء تسمى بـ "المعجزات"، ووضعوا لها قيوداً وشروطاً لم تأتِ في القرآن ولا في السنة، كقولهم: "أن تتحدى بها قوماً"، "وأن يعجزوا عنها"، "وأن تكون مما برع فيه القوم".

الصواب أن آيات الأنبياء سميت في القرآن (آيات، حججاً، براهين، سلطاناً، بصائر). قال تعالى:
\begin{ayah}
وَمَا كَانَ لِرَسُولٍ أَن يَأْتِيَ بِآيَةٍ إِلَّا بِإِذْنِ اللَّهِ
\end{ayah}
وقال النبي صلى الله عليه وسلم:
\begin{hadith}
ما من الأنبياء من نبي إلا قد أُعطي من الآيات ما مثله آمن عليه البشر، وإنما كان الذي أوتيت وحياً أوحاه الله إليّ.
\end{hadith}

الغاية من هذه الآيات ليس مجرد "التعجيز" أو "التحدي الابتدائي"، بل الغاية هي الهداية وبيان الحق؛ كما قال تعالى: \begin{ayah}سَنُرِيهِمْ آيَاتِنَا فِي الْآفَاقِ وَفِي أَنفُسِهِمْ حَتَّىٰ يَتَبَيَّنَ لَهُمْ أَنَّهُ الْحَقُّ\end{ayah}. فلما ادعوا كذباً وزوراً \begin{ayah}لَوْ نَشَاءُ لَقُلْنَا مِثْلَ هَٰذَا\end{ayah}، تحداهم الله تبارك وتعالى أن يأتوا بسورة من مثله، فظهر عجزهم كبرهان على أنه من عند الله.

\begin{fawaid}[نقد مصطلح المعجزة]
مصطلح "المعجزة" اصطلاح حادث أوجده المتكلمون. وفيه قصور دلالي؛ لأن العجز وحده ليس دليلاً على أن الشيء من عند الله (فالسحر مثلاً يعجز عنه عموم الناس وليس من الله). اللفظ الشرعي الأكمل هو "الآية"، فهو يتضمن معنى العجز وزيادةً عليه كونه يحمل العلم والهدى والرحمة والبرهان على صدق المرسِل.
\end{fawaid}

\separator

\section{نقد مذاهب المتكلمين في إعجاز القرآن}
ذهب المتكلمون مذاهب شتى في تحديد وجه إعجاز القرآن، وكثير من كلامهم مأخوذ من الفلسفات ومخالف للسان العرب. وفي هذا الموضع عرض الشيخ بعض هذه المذاهب، ثم تعقبها، وتوقف خصوصاً عند ما يفهم من بعض كلام \rawi{أحمد شاكر} رحمه الله في حصر جهة التحدي في البيان وحده.

\subsection{أولاً: بطلان القول بالصَّرْفَة}
ذهب كثير من المعتزلة إلى أن الإعجاز في "الصَّرْفَة"، وهي كلمة أحدثوها تعني أن الله سبحانه صرف الناس عن أن يأتوا بمثل القرآن، وكأنهم في الأصل كانوا قادرين على الإتيان بمثل كلام الله ولكن الله قهرهم وصرفهم! وهذا غلو في الإنسان مأخوذ من الفلاسفة الذين يجعلون البشر آلهة مقهورة.

\subsection{ثانياً: قصور القول بحصر الإعجاز في "النظم والبلاغة"}
ذهب \rawi{عبد القاهر الجرجاني} وجماعة من الأشاعرة إلى حصر إعجاز القرآن في اعتباره إعجازاً بلاغياً أو في نظمه ووزنه. وهذا حق في ذاته، ولكنه يعاني من مشكلتين:
 \begin{enumerate}
    \item أنهم جعلوه "كل" إعجاز القرآن، وهذا قصور؛ فالقرآن معجز في نظمه وبيانه، وفي علمه وإخباره بالغيب، وفي تشريعه وأحكامه.
    \item أنهم أتعبوا أنفسهم في محاولة "التعبير عن الفرق" بين كلام الله وكلام البشر. والقرآن كالحياة، ندرك الفرق بين الحي والميت، ولكننا لا نحيط علما بحقيقة الحياة. فكذلك القرآن يُعلم ببديهة العربي أنه من عند الله وليس من كلام البشر، ولكن البشر يعجزون عن التعبير عن هذا الفرق لأنهم لا يحيطون بكلام الله علماً.
\end{enumerate}

ومحل العناية هنا أن الشيخ لم يكتف بإثبات الإعجاز البياني، بل نبه على أن في القرآن جهات أخرى ظاهرة من الإعجاز، مثل ما فيه من أخبار الغيب، وما فيه من التشريع والهدى. فحصر التحدي في مجرد النظم والبيان وحده تضييق لما وسعه القرآن نفسه.

\begin{fawaid}[لا يحصر إعجاز القرآن في جهة واحدة]
من أقوى تعقيبات الشيخ في هذا اللقاء أن القرآن ليس معجزاً من جهة البيان وحدها، بل من جهة ما تضمنه أيضاً من العلم والهدى والصدق في الأخبار والعدل في الأحكام. فالإعجاز البياني أصل ظاهر، لكن قصر الإعجاز عليه وحده يفضي إلى نقص في تصور حقيقة القرآن.
\end{fawaid}

\separator

\section{الألفاظ المشتركة بين العرب والعجم}
دخل \rawi{الطبري} رحمه الله في معالجة ظاهرة لغوية هامة: وهي وجود كلمات في القرآن يستعملها العرب ويستعملها غيرهم من الأعاجم (كالفرس والروم والحبشة)، مثل: (القسطاس، سجيل، مشكاة).

كيف نفسر هذه الظاهرة في ضوء أن القرآن نزل \begin{ayah}بِلِسَانٍ عَرَبِيٍّ مُّبِينٍ\end{ayah}؟
ذكر \rawi{الطبري} الاحتمالات:
\begin{itemize}
    \item هل أخذها العرب من العجم فعربوها؟
    \item أم أخذها العجم من العرب؟
    \item أم هي من توافق اللغات والألسنة؟
\end{itemize}

مال \rawi{الطبري} إلى القول بأنها من باب "اتفاق الألسنة"، وأنها ألفاظ اشتركت فيها الأمم في اللفظ والمعنى معاً. وقال: «الصواب في ذلك عندنا أن يسمى عربياً أعجمياً، أو حبشياً عربياً... وليس غير ذلك من كلام كل أمة منهما بأولى أن يكون إليها منسوباً».

وقد رد على الآثار المروية عن \rawi{ابن عباس} وغيره (في أن بعض الكلمات حبشية أو فارسية) بأن الصحابي حين نَسَب اللفظ للحبشة، فهو يُثبت أن الحبشة تنطق به، ولم ينفِ أنه عربي من كلام العرب أيضاً.

\begin{fawaid}[التفريق بين الظاهرة وتفسيرها]
هناك فرق دقيق بين "الظاهرة" و"تفسير الظاهرة". الظاهرة هنا: وجود كلمات متطابقة في اللغات. هذا متفق عليه. أما "كيف حدث هذا التطابق؟" فهذا هو التفسير الذي يحتاج إلى دليل. وقد اختار \rawi{الطبري} أنه توارد ألسنة، بينما اختار آخرون أن العرب استعارتها من جيرانها وعربتها قبل نزول القرآن، فصارت بهذا التعريب عربية فصيحة نزل بها القرآن، وهو القول الأقوى.
\end{fawaid}

\begin{fawaid}[قاعدة في اختلاف التضاد والتنوع]
إذا أمكن اجتماع الوصفين بجهتين مختلفتين فهو "اختلاف تنوع" ولا يصح ضرب الأقوال ببعضها. فإذا قيل للكلمة "عربية" (أي تنطق بها العرب) وقيل "حبشية" (أي تنطق بها الحبشة)، فلا تضاد. الإثبات لا يدل على النفي فيما يجوز اجتماعه من المعاني، كما يقال لأرض أنها "سهلية وجبلية" في ذات الوقت لاشتمالها على الوصفين.
\end{fawaid}

\separator

\section{خطورة التسليم للمخالفين في الأصول}
أثناء استعراض الدكتور \rawi{مساعد الطيار} لبعض الأساليب العربية كالمشاكلة اللفظية (مثل: \begin{ayah}وَجَزَاءُ سَيِّئَةٍ سَيِّئَةٌ مِّثْلُهَا\end{ayah})، حذّر من أن ننفي وجود هذا الأسلوب في القرآن بالكلية فقط لأن بعض المنحرفين عقديّاً استخدموه لتأويل آيات الصفات (مثل: \begin{ayah}وَيَمْكُرُونَ وَيَمْكُرُ اللَّهُ\end{ayah}).

وتعليقاً على ذلك نقول: من الخطأ المنهجي الجسيم أن توافق المخالف على "القاعدة الكلية" ثم ترفض تطبيقها في جزئية معينة! 
فإذا سلمت للمتكلمين بأن القرآن فيه "مشاكلة لفظية" أو مجاز مطلق، ثم جئت عند آيات الصفات وقلت: (إلا في الصفات لا يجوز المجاز)، سيقولون لك: ما دليلك على الاستثناء؟

والصواب عند الشيخ هو ما فعله \rawi{ابن تيمية} رحمه الله؛ وهو هدم هذه الأصول والقواعد المبتدعة من أساسها، لا مجرد التسليم لها ثم محاولة منع بعض نتائجها. فالفروع الباطلة ناتجة عن أصول خاطئة يجب اجتثاثها.

\separator

بهذا نكون قد أتممنا الجزء الثاني من مقدمة \rawi{الطبري}، وغداً إن شاء الله ننهي المقدمة ونبدأ في التفسير العملي لسورة الفاتحة. والسلام عليكم ورحمة الله وبركاته.
